% ======================================================================
% 01-intro.tex
% ======================================================================
\section{Introduction}
\label{sec:intro}

Three numbers illustrate why this article is about software rather than
physics.

The first is a Fourier neural operator (FNO) with relative $L^2$ error
$0.973$. It had been recorded as a weak result. It was not weak: because the
metric is normalized by the target norm, a predictor of all zeros scores exactly
$1.0$, so the operator was \emph{worse than predicting nothing}. The input was
two spatially constant parameter channels; a pointwise lift followed by spectral
layers that multiply Fourier coefficients can only produce a constant output.
The architecture was structurally incapable of representing a front, and no
amount of training could fix it.

The second is a physics-informed neural network (PINN) reported at
$1.35\times10^{-2}$ full-grid error and $8.0\times10^{-3}$ in the interior. It
was the best of five seeds. At fixed hyperparameters the same configuration
spans $1.17\times10^{-2}$ to $8.14\times10^{-2}$---a factor of seven---and the
median interior error is $1.42\times10^{-2}$, not below the $10^{-2}$ target the
result was said to meet.

The third is a claimed $5\times$ GPU speedup that had never been measured. When
the launch count was finally instrumented, one evaluation issued $4{,}022$
batched matrix--vector products, $1.7\times$ the number the optimization plan
had assumed. The speedup may still be real; the point is that it was being
\emph{planned} as if it were a result.

None of these is a modelling failure. Each is a failure of the process that
decides what counts as evidence: a metric without a baseline, a run without
repetition, a claim without an instrument. The response could have been a
refactor. Instead, this project accumulated \emph{gates}: a trivial-predictor
guard that makes an FNO worse than zeros fail fast; a seed-sweep harness that
refuses to compare unpaired runs; a backend-parity suite that treats a numerical
disagreement as a bug; an artifact-first remote runner that keeps the results of
a run that dies. The gates are the reusable code.

\paragraph{Why AI matters here.}
Writing research software is usually a choice between doing the experiment and
doing the software hygiene, and the experiment wins. The standard remedies---
version control, tests, continuous integration, documented environments---are
well known \citep{wilson2014best,sandve2013ten}, but they are labour, and the
labour is repetitive rather than creative. An agent that can read the repository,
run the test, read the failure, edit the file, rerun, and write down what
happened changes the marginal cost of that labour. Recent work on agents that
act through tools and on benchmarks built from real repository issues makes the
same point from the other direction: the hard part is not generating a patch,
it is closing the loop against an executable check
\citep{yao2023react,jimenez2024swebench}.

This article is an experience report from one project in which that loop was
closed with the Pi coding agent\footnote{\url{https://pi.dev/}} running DeepSeek
models.\footnote{\url{https://www.deepseek.com/}} The agent performed most of
the edits, test runs, benchmark harnesses and documentation updates. The human
chose the gates, rejected the tempting numbers, and wrote the final judgements.
We are explicit about this division because it is the article's main claim.

\paragraph{Thesis.}
Reusability is not a property of a directory layout or a test count. It is the
accumulation of falsifiable gates, written memory, and one reference
implementation---and an AI agent made that accumulation cheap enough to do on
every failure instead of once at the end.

\paragraph{Claims.}
The article argues three claims, each with evidence from the repository:
\begin{enumerate}
  \item \textbf{Reusability is measurable.} The same policy reproduces on a
        third machine, the NumPy and torch backends agree to $10^{-10}$ and
        select the same policy, and a certificate on 400 unseen scenarios
        survives a redraw.
  \item \textbf{The reusable artifacts came from a repeatable loop.} Every
        mechanism in \S\ref{sec:mechanisms} and every guard in
        \S\ref{sec:failures} is traced to a specific defect and a commit; none
        arrived from a big-bang rewrite.
  \item \textbf{AI's contribution was verification throughput and memory, not
        code volume.} The failure catalogue includes cases where the agent was
        confidently wrong, and the human protocol that caught them is part of
        the method.
\end{enumerate}

\paragraph{Contributions.}
(i) A concrete account of five reusability mechanisms in a computational
hydrology stack: reference/batched parity, configuration manifests, gates,
statistical harnesses, and artifact-first remote execution. (ii) A failure
catalogue in which each defect became a permanent guard, with the commit and
the test that pins it. (iii) A set of results that survive independent
re-derivation, including one that fails its own in-run certificate. (iv) A
practical account of what an AI agent did well and where it failed.

\paragraph{Roadmap.}
\S\ref{sec:case} fixes the scientific setting. \S\ref{sec:method} describes the
AI-assisted loop and the artifacts that serve as memory. \S\ref{sec:mechanisms}
presents the five mechanisms. \S\ref{sec:failures} is the failure catalogue.
\S\ref{sec:results} reports the results that survived, with their caveats.
\S\ref{sec:discussion} discusses the collaboration and its threats to validity,
and \S\ref{sec:conclusion} concludes. Appendices collect gate definitions,
repository layout and full tables.
